\section{Revisione di sicurezza}\label{sec:revisione}

Completate le funzioni principali, il prototipo è stato sottoposto a una revisione di sicurezza complessiva, per cercare i problemi che i test scritti insieme al codice non avevano rilevato.

\subsection{Metodo}
La revisione ha riguardato l'intero repository: backend, frontend, contratti, configurazione di Docker e MongoDB, segreti, dipendenze e storia di Git. Come riferimento si è usato lo standard di verifica OWASP ASVS~5.0\cite{asvs} e, per i contratti, le vulnerabilità note degli smart contract\cite{atzei-survey}. Le verifiche sono state di cinque tipi:
\begin{itemize}
  \item rilettura del codice, controllo per controllo;
  \item \textbf{riproduzione} di ogni sospetto con una prova contro il codice reale: un sospetto non riprodotto non è stato classificato come vulnerabilità;
  \item analisi delle dipendenze con \texttt{npm audit} e \texttt{pip-audit}, e verifica che ogni componente sia ancora supportato dal produttore;
  \item ricerca di segreti in tutta la storia del repository con \texttt{gitleaks} e con ricerche mirate;
  \item analisi statica del codice TypeScript con CodeQL.
\end{itemize}
Per ogni problema confermato si è scritto prima un test che lo riproduce e fallisce, poi la correzione; il test resta nella suite, così il problema non può ripresentarsi.

\subsection{Risultati}
La tabella~\ref{tab:revisione} riassume i problemi confermati. La gravità tiene conto di quanto il problema è facile da sfruttare e di cosa ne ottiene un attaccante.

\begin{table}[H]
\centering\footnotesize
\renewcommand{\arraystretch}{1.2}
\begin{tabular}{@{}p{.03\textwidth}p{.31\textwidth}p{.08\textwidth}p{.37\textwidth}p{.09\textwidth}@{}}
\toprule
\# & Problema & Gravità & Correzione & Stato\\
\midrule
1 & Chiunque poteva bloccare il login di un utente sostituendone la challenge & Alta & Challenge indicizzate per valore e lette dai dati firmati & Corretto\\
2 & Errori interni inviati al client con stack trace e messaggi del database & Alta & Gestore finale degli errori con messaggi generici & Corretto\\
3 & Registro degli accessi in chiaro nel database & Alta & Voci cifrate e legate alla loro impronta on-chain & Corretto\\
4 & Le opzioni di login rivelavano se un utente esisteva & Media & Opzioni uguali per tutti, passkey individuabili & Corretto\\
5 & Nessun limite di frequenza; ogni richiesta di cartella produce una transazione & Media & 60 richieste al minuto per client e per utente & Corretto\\
6 & Nessuna intestazione di sicurezza; dati sanitari conservabili in cache & Media & \texttt{no-store}, CSP, \texttt{nosniff}, protezione dai frame, HSTS in produzione & Corretto\\
7 & Contatore anti-riuso aggiornato senza controllo di concorrenza & Media & Compare-and-set sul valore atteso & Corretto\\
8 & Nessuna revoca di passkey o sessioni, né durata massima della sessione & Media & Revoca, uscita da tutti i dispositivi, durata massima di 8 ore & Corretto\\
9 & Un'unica chiave di amministrazione dei contratti, cedibile in un attimo & Media & Passaggio in due fasi con attesa, pausa con ruolo separato & Corretto\\
10 & Quote del backup delle chiavi pubblicabili per errore & Media & Escluse da Git; lo script rifiuta le cartelle del progetto & Corretto\\
11 & Metadati collegabili sulla blockchain pubblica & Media & Pseudonimi; serve una scelta architetturale & Residuo\\
12 & Chiavi di prova di Hardhat utilizzabili su una rete reale & Bassa & Avvio rifiutato fuori dalla rete locale & Corretto\\
13 & Parametri di Argon2id letti dal file senza limiti & Bassa & Parametri accettati solo entro limiti & Corretto\\
14 & Tentativi d'emergenza respinti non registrati & Bassa & Conservati nel registro cifrato & Corretto\\
15 & Derivazione del KEM ibrido senza le chiavi pubbliche del destinatario & Bassa & Derivazione conforme a SP~800-227 & Corretto\\
16 & Immagine Docker indicata per nome e non per impronta & Bassa & Immagine fissata con il \emph{digest} & Corretto\\
17 & Nomi utente nei log del server & Bassa & Nessun dettaglio clinico nei log & Residuo\\
18 & Vulnerabilità in una dipendenza di sviluppo dei contratti & Bassa & Vedi sotto & Residuo\\
\bottomrule
\end{tabular}
\caption{Problemi emersi dalla revisione di sicurezza.}\label{tab:revisione}
\end{table}

\paragraph{Dipendenze e segreti.} Le dipendenze di backend, frontend e notebook, cioè il codice che va in esecuzione, non hanno vulnerabilità note. Tra gli strumenti di sviluppo dei contratti resta una segnalazione su \texttt{zod}, la libreria con cui Hardhat controlla il proprio file di configurazione: la versione corretta non è compatibile con Hardhat. Il difetto riguarda il consumo di memoria su input forniti dallo sviluppatore stesso e non tocca il servizio; va ricontrollato a ogni aggiornamento di Hardhat. La ricerca di segreti in tutti i commit ha trovato solo le chiavi di prova pubbliche di Hardhat, valide sulla sola rete locale.

\paragraph{Analisi statica.} CodeQL ha segnalato come prive di limite di frequenza cinque rotte di login, registrazione e firma. Le segnalazioni sono state verificate una per una e archiviate come falsi positivi: il limite è applicato a tutte le rotte del gruppo da un middleware del progetto, che CodeQL non riconosce, e un test ne dimostra l'effetto (risposta 429 oltre la soglia). L'analisi statica segnala schemi senza conoscere l'intenzione, per questo ogni segnalazione va verificata e non accettata o ignorata in blocco.

\subsection{Riepilogo dei test}
La tabella~\ref{tab:test-totali} riporta i test automatici del progetto, tutti superati. I flussi visibili all'utente (registrazione, accesso, gestione delle passkey, consultazione e firma dei referti, consensi dal wallet, verifica del registro sulla catena) sono stati provati anche in un browser reale, con un autenticatore virtuale e con un autenticatore di piattaforma.

\begin{table}[H]
\centering\footnotesize
\renewcommand{\arraystretch}{1.2}
\begin{tabular}{@{}p{.36\textwidth}p{.32\textwidth}r@{}}
\toprule
Gruppo & Comando & Test\\
\midrule
Backend, unitari & \texttt{npm test} (cartella \texttt{backend}) & 127\\
Backend su MongoDB reale & \texttt{npm run test:mongo} & 20\\
Backend sulla blockchain locale & \texttt{npm run test:chain} & 8\\
Smart contract & \texttt{npm test} (cartella \texttt{blockchain}) & 36\\
\midrule
Totale & & 191\\
\bottomrule
\end{tabular}
\caption{Test automatici del progetto.}\label{tab:test-totali}
\end{table}

\subsection{Rischi residui}\label{sec:residui}
\begin{itemize}
  \item \textbf{Metadati sulla catena pubblica.} Anche con indirizzi pseudonimi, gli eventi collegano paziente, medico, finalità e orario: se è noto l'indirizzo di un oncologo, si può dedurre chi è in cura da lui. Le mitigazioni sono uno pseudonimo per ogni rapporto paziente--medico, prove a conoscenza zero (sezione~\ref{sec:zkp}) o una rete \emph{permissioned} (sezione~\ref{sec:fabric}). È la ragione principale per evitare la rete pubblica in un'adozione reale.
  \item \textbf{Backend compromesso.} Chi controlla il backend può leggere i dati che il backend decifra senza registrarne l'accesso, ma non può concedere consensi, riscrivere l'audit o far risultare autorizzato un accesso senza consenso.
  \item \textbf{Operazioni sensibili senza nuova verifica.} La revoca di una passkey richiede solo una sessione valida; una nuova verifica con la passkey (\emph{step-up}) proteggerebbe una postazione lasciata aperta.
  \item \textbf{Produzione.} Servirebbero TLS verso il database, cifratura dei volumi, HTTPS verso i client, log con conservazione limitata e un controllo periodico che ogni componente sia supportato; il repository ha già attivi gli avvisi sulle dipendenze (Dependabot), il blocco dei \emph{push} con segreti e l'analisi statica.
\end{itemize}
Gli altri rischi residui, specifici di un componente, sono discussi nelle sezioni~\ref{sec:fido2}, \ref{sec:blockchain} e~\ref{sec:dati}.
